\section*{Chapitre \ref{ch:s1:alternative-data-strategies} --- Stratégies fondées sur des données alternatives}

\begin{solution}{exo:s1:alternative-data-strategies:1}
Avec un bruit deux fois la dispersion de la surprise, la corrélation vaut $1/\sqrt{1 + 2^2} = 0.45$ et la meilleure pente $1/(1 + 2^2) = 0.2$~: le \omterm{def:s1:alternative-data-strategies:nowcast}{nowcast} réduit la lecture
de quatre cinquièmes. Les valeurs ajustées du chapitre sont 0.46 et 0.21.
\end{solution}

\begin{solution}{exo:s1:alternative-data-strategies:2}
Le \omterm{def:s1:alternative-data-strategies:nowcast}{nowcast} vaut $0.2 \times 2 = 0.4$ écart-type de surprise, et la réaction attendue $3 \times 1.6\% \times 0.4 = 1.92\%$. La moitié du marché réagit
à hauteur de $0.5 \times 1.92~\% = 0.96~\%$ à la livraison, l'autre moitié étant laissée au rapport.
\end{solution}

\begin{solution}{exo:s1:alternative-data-strategies:3}
La décroissance après publication suit la diffusion d'une idée que chacun peut recopier sans frais~; l'\omterm{def:s1:alternative-data-strategies:decay}{érosion des données} suit le nombre d'acheteurs d'un jeu de données dont l'accès est payant
et peut être vendu en exclusivité. Son rythme peut être ralenti par l'exclusivité et la rapidité, et il peut être estimé à partir des ventes du jeu de données plutôt que du calendrier.
\end{solution}

\begin{solution}{exo:s1:alternative-data-strategies:4}
Le portefeuille tourne 36 fois son exposition brute par an, à sens unique, à dix points de base par unité négociée~; lorsque l'essentiel de la réaction a lieu à la livraison, le rendement brut passe sous les
coûts avant que la part des acteurs disposant des données n'atteigne 1~: à 0.9, le ratio de Sharpe brut vaut 0.15 et le ratio de Sharpe net $-0.43$.
\end{solution}

\begin{solution}{exo:s1:alternative-data-strategies:5}
Avec les données un jour plus tôt, la société est positionnée avant que les autres n'opèrent sur elles, si bien qu'elle gagne leur mouvement à la livraison ainsi que le reste lors du rapport~:
le total de la réaction reste inchangé, et la société garde sa position pendant toute la réaction, quelle que soit la part du marché qui suit.
\end{solution}

\begin{solution}{exo:s1:alternative-data-strategies:6}
Mesurer la variation du cours les jours de livraison (ou dans les jours précédant les rapports) par unité de \omterm{def:s1:alternative-data-strategies:nowcast}{nowcast}~: si personne d'autre n'utilise les données, elle est nulle~; à mesure que l'usage
se développe, elle tend vers la réaction attendue dans son ensemble. Le rapport entre le mouvement d'avant rapport et le mouvement total permet d'estimer la part.
\end{solution}

\begin{solution}{exo:s1:alternative-data-strategies:7}
Avec un bruit trois fois la dispersion de la surprise (une corrélation de 0.34), le portefeuille à demi-diffusion dégage un ratio de Sharpe de 1.35 et 17.6~\% par an, contre 1.86 et 24.2~\% avec
un bruit deux fois cette dispersion~; sans diffusion, 3.20 contre 4.41. Un panel plus bruité prédit moins de chaque surprise, si bien que chaque position porte une part moindre de la réaction
et supporte les mêmes coûts.
\end{solution}

\begin{solution}{exo:s1:alternative-data-strategies:8}
L'historique d'un fournisseur peut avoir été complété rétroactivement (reconstitué après coup et ajusté aux résultats), et il ne dit pas combien d'autres clients possèdent désormais les
données. Exiger des livraisons point-in-time, tester les données en conditions réelles pendant une période, et estimer leur diffusion actuelle à partir des mouvements d'avant rapport.
\end{solution}

\begin{solution}{pb:s1:alternative-data-strategies:1}
\begin{enumerate}
\item Une estimation d'une grandeur publiée ultérieurement à partir de données antérieures~; une stratégie de négociation fondée sur de tels \omterm{def:s1:alternative-data-strategies:nowcast}{nowcasts}, issus de données non traditionnelles.
\item Rattacher les données aux sociétés, estimer par \omterm{def:s1:alternative-data-strategies:nowcast}{nowcast} la grandeur publiée, la traduire en réaction attendue, et tenir la position sur la fenêtre pendant laquelle le marché l'apprend.
\item Un dixième des sociétés, des lectures bruitées à hauteur de deux fois la dispersion de la surprise (corrélation 0.46), livrées dix jours à l'avance~; pente 0.21 estimée sur les deux premières
années.
\item Positions proportionnelles au \omterm{def:s1:alternative-data-strategies:nowcast}{nowcast}, de la clôture de la livraison à la clôture de l'annonce, exposition brute de 1.
\item Les données satellitaires permettent à des investisseurs sophistiqués de cibler les mauvais rapports, avec davantage de ventes à découvert informées et une liquidité moindre~: plus d'asymétrie, et non
des prix immédiatement meilleurs.
\item Les notes d'employeurs en amélioration surperforment celles en baisse et annoncent les surprises de résultats.
\item Les cours ignorent les nouvelles concernant les principaux clients~; une stratégie fondée sur les fournisseurs a rapporté plus de 150 points de base par mois.
\item Les données alternatives ont accru le caractère informatif des prix en réduisant les coûts d'information.
\item La perte de valeur à mesure que davantage d'investisseurs négocient les mêmes données plus tôt~; une part $\phi$ déplace $\phi$ de la réaction attendue à la livraison.
\item 4.41, 3.20, 1.86, 0.45 et $-0.43$ pour des parts de 0, 0.25, 0.5, 0.75 et 0.9.
\item Chaque acheteur déplace à la livraison une part égale de la réaction, hors de portée du portefeuille.
\item Il maintient le ratio de Sharpe autour de 4.3--4.4 quel que soit le niveau de diffusion.
\item \textbf{Résultat nommé.} Le ratio de Sharpe net de coûts du portefeuille de \omterm{def:s1:alternative-data-strategies:nowcast}{nowcasts} passe de 4.41 lorsque personne d'autre ne dispose du panel à 1.86 lorsque la moitié du marché le détient, et
$-0.43$ à neuf dixièmes, tandis qu'une société qui reçoit le panel un jour plus tôt conserve 4.28 à 4.43.
\item Suivre le mouvement d'avant rapport par unité de \omterm{def:s1:alternative-data-strategies:nowcast}{nowcast} et la réaction réalisée du portefeuille par transaction.
\item Livraisons tardives ou modifiées, pertes de couverture, révisions, défaillances du fournisseur et conditions de licence.
\item Exiger un historique point-in-time, vérifier l'absence de remplissage rétroactif, et le tester en conditions réelles.
\item Ce qu'il ajoute aux rendements de la société après coûts sur sa durée de vie attendue, compte tenu de sa diffusion.
\item La stratégie de \omterm{def:s1:alternative-data-strategies:nowcast}{nowcast} du chiffre d'affaires, dont les données sont vendues à de nombreux acheteurs.
\item Elles opèrent sur les mêmes rapports avant leur publication~; le chapitre~7 les traite après.
\item Le temps~: la connaissance d'une grandeur avant que le marché ne la possède.
\end{enumerate}
\end{solution}

\begin{solution}{iq:s1:alternative-data-strategies:1}
Couverture et historique (point-in-time~?), rattachement aux sociétés, corrélation avec la grandeur publiée sur des données que le fournisseur n'a pas pu ajuster, le
gain incrémental de l'IC par rapport aux signaux existants de la société, le décalage, le coût, qui d'autre l'achète, et les conditions juridiques.
\end{solution}

\begin{solution}{iq:s1:alternative-data-strategies:2}
Historique reconstitué après coup, souvent ajusté à des résultats connus, si bien que le backtest paraît meilleur qu'un usage en conditions réelles. On le détecte en comparant l'ajustement avant et
après la date de lancement du fournisseur, par les sauts de couverture et les ruptures de niveau au lancement, et en demandant les horodatages de livraison (livre~7, chapitre~12).
\end{solution}

\begin{solution}{iq:s1:alternative-data-strategies:3}
Davantage d'acheteurs (diffusion), un changement de panel (fournisseur, couverture, méthodologie), un changement dans ce que déclarent les sociétés, ou un encombrement (crowding) dû à des jeux de données similaires.
Vérifier le mouvement d'avant rapport par unité de signal et les statistiques propres au panel.
\end{solution}

\begin{solution}{iq:s1:alternative-data-strategies:4}
Ingérer et valider la livraison (schéma, couverture, plages de valeurs par rapport à l'historique), rattacher aux identifiants des sociétés point-in-time, calculer les \omterm{def:s1:alternative-data-strategies:nowcast}{nowcasts} avec le modèle
en vigueur, les transmettre à l'optimiseur avec les autres signaux, générer les ordres pour la clôture, et tout consigner pour permettre la reproduction.
\end{solution}

\begin{solution}{iq:s1:alternative-data-strategies:5}
L'information privilégiée (material non-public information, MNPI), c'est-à-dire des données obtenues en violation d'un devoir, les données personnelles et le droit de la vie privée, les conditions de licence et l'exclusivité, ainsi que les conditions d'utilisation des
services de web scraping~; chaque jeu de données exige un contrôle de conformité sur la façon dont il a été collecté.
\end{solution}

\begin{solution}{iq:s1:alternative-data-strategies:6}
Avec une réaction attendue $J$ par unité de \omterm{def:s1:alternative-data-strategies:nowcast}{nowcast}, un trader à la livraison gagne $(1 - \phi)J$ par unité, puisque $\phi J$ se déplace lorsque les autres négocient. Un trader disposant des
données $k$ jours plus tôt, qui conserve sa position jusqu'au bout, gagne la totalité de $J$ (mouvement des autres compris), moins toute dérive des jours supplémentaires~; la valeur de la rapidité est $\phi J$.
\end{solution}
